\documentclass[a4paper,12pt]{article}

% Set margins
\usepackage[hmargin=1cm, vmargin=0.5cm]{geometry}

\frenchspacing

% Language packages
\usepackage[utf8]{inputenc}
\usepackage[T1]{fontenc}
\usepackage[magyar]{babel}

% AMS
\usepackage{amssymb,amsmath}

% Graphic packages
\usepackage{graphicx}

% Colors
\usepackage{color}
\usepackage[usenames,dvipsnames]{xcolor}

% Enumeration
\usepackage{enumitem}

% Links
\usepackage{hyperref}

\newcommand{\code}[1]{\texttt{#1}}

\pagestyle{empty}

\begin{document}

\begin{center}
	
	\
	
	\vskip 1cm
	
	{\huge \textbf{Záróvizsga tételek}}
	
	\medskip

	{\large \textit{Programtervező informatikus BSc}}
\end{center}

\vskip 5mm

\renewcommand{\arraystretch}{1.2}

\begin{tabular}{|lp{16.5cm}|}

\hline

\textbf{1.} &
Adatok, adattípusok, adatműveletek és adatstruktúrák. Szám adattípus. Számrendszerek, konverziók. A logikai, halmaz, karakter, sztring absztrakt adattípusok és realizációjuk. A tömb (vektor, mátrix), rekord absztrakt adattípusok.
\\

\textbf{2.} &
Az algoritmus. Iteratív és rekurzív algoritmus. A számítógépes memória. Adat és program. Verem és procedúra. Az algoritmus lejegyzése. A folyamatábra és a pszeudokód. Elemi algoritmusok.
\\

\textbf{3.} &
Strukturált programozás. Programgráf, valódi program, vezérlőgráf lebontása, strukturált program és annak formulája. Strukturált programgráf kialakítása, struktogram. Ciklikus bonyolultság és egyéb bonyolultsági tételek.
\\

\hline

\textbf{4.} &
Számelméleti algoritmusok. Legnagyobb közös osztó, euklideszi és kibővített euklideszi algoritmus, lineáris kongruencia egyenletek. Multiplikatív inverz, moduláris hatványozás, Fermat prímteszt. RSA.
\\

\textbf{5.} &
Rendezések: Beszúró rendezés. Az "Oszd meg és uralkodj!" elv. Összefésülő rendezés. Gyors rendezés. Buborék rendezés. Shell rendezés. Minimum kiválasztásos rendezés. Négyzetes rendezés. Lineáris idejű rendezések: leszámláló rendezés, számjegyes rendezés. Időelemzéseik. Az összehasonlító rendezések időtétele.
\\

\textbf{6.} &
Gráf algoritmusok. Szélességi keresés. Mélységi keresés. Topológikus rendezés. Optimum feladatok fákon (minimális feszítőfák, a Kruskal és Prim algoritmus). Legrövidebb utak meghatározása (Dijkstra algoritmus, Bellman-Ford algoritmus. Floyd-Warshall algoritmus).
\\

\hline

\textbf{7.} &
A párhuzamos algoritmusok alapfogalmai, PRAM modell, hatékonysági mértékek (munka, költség, gyorsítás, hatékonyság). Párhuzamos végrehajtás ábrázolási módjai. Komplexitás becslése, párhuzamos komplexitás. Amdahl törvénye. Pipeline párhuzamosítás, Lost update probléma bemutatása és megoldása.
\\

\textbf{8.} &
Aggregáció, redukció, prefixszámítás párhuzamosítási lehetőségei. CREW\_PREFIX, EREW\_PREFIX és OPTIMAL\_PREFIX algoritmusok bemutatása. Összefésülő rendezés párhuzamosítási módja.
\\

\textbf{9.} &
A busz, gyűrű, rács, tórusz, hiperkocka és fa topológiák bemutatása, csomópontok közötti távolságok jellemzése. A CSP, Publish-Subscribe és az Aktor modell. Az aszinkron programozás elemei: \code{async}, \code{await} kulcsszavak használata, coroutine-ok.
\\

\hline

\textbf{10.} &
A Turing gép fogalma, működése. A RAM-gép. Boole-függvények és logikai hálózatok.
\\

\textbf{11.} &
Algoritmikus eldönthetőség. Church-tézis. Rekurzív és rekurzívan felsorolható nyelvek, rekurzív illetve parciálisan rekurzív függvények. Nevezetes nyelvek (Az R, Re, coR, coRE nyelvosztályok és ezek kapcsolata) és bonyolultságuk. Algoritmikusan eldönthetetlen problémák. Polinomiális idejű algoritmusok.
\\

\textbf{12.} &
Nemdeterminisztikus Turing gépek, Az NP és a coNP nyelvosztály. Példák NP-beli nyelvekre.  A tanú-tétel. Nemdeterminisztikus algoritmusok bonyolultsága. NP-teljesség, Cook-tétel.  Néhány NP-teljes probléma, Cook-Levin tétel.
\\

\hline

\end{tabular}

\begin{tabular}{|lp{16.5cm}|}

\hline

\textbf{13.} &
A Java programozási nyelv története, alapvető jellemzői. Egy Java program szerkezete (csomagok és fordítási egységek). Az objektum orientált programtervezés alapelveinek felsorolása. Objektumok tárolása és élettartama. Szemétgyűjtő mechanizmus. Kivételkezelés.
\\

\textbf{14.} &
Az egységbezárás és az információrejtés objektum orientált alapelvek megvalósítása. Osztálydefiníció és példányosítás. Osztályszintű és példányszintű tagok, konstansok. Hozzáférési kategóriák és használatuk. Hivatkozás az osztály tagjaira.
\\

\textbf{15.} &
Az öröklődés és a polimorfizmus objektum orientált alapelvek megvalósítása. Öröklődés szabályai, metódusok túlterhelése és felüldefiniálása. A referencia statikus és dinamikus típusa. Absztrakt osztály és interfész szerepe, összehasonlítása.
\\

\hline

\textbf{16.} &
Az operációs rendszerbeli folyamat (processz) fogalom. Folyamat kontextus. A fonál (szál) fogalom. A processz állapotok, állapotváltások, processz futási módok. Taszk és fonál állapotok, állapotátmenetek.
\\

\textbf{17.} &
A memória menedzselés feladatai. Memória, mint erőforrás. Címleképzés fajták. Lapozós virtuális memória menedzselés működése. Allokálás, nyilvántartás, címleképzés. Laphiba fogalma, kezelése, kilapozási stratégiák. Előnyök, hátrányok.
\\

\textbf{18.} &
Fájlrendszer megvalósítási feladatok. Jegyzékszerkezetek. Szabad blokk menedzselési lehetőségek. Fájl attribútum rögzítési lehetőségek (ezen belül fájl testet képező blokkok rögzítési lehetőségei).
\\

\hline

\textbf{19.} &
Relációs adatmodell, relációs struktúra és integritási feltételek. ER modell elemei jelentése és jelölése. Az ER modell konverziója relációs modellre. Relációs adatmodell műveleti része, relációs algebra.
\\

\textbf{20.} &
Az SQL szabvány relációs adatbázis kezelő nyelv bemutatása, a DDL (objektum létrehozás, megszüntetés és szerkezet módosítás) , DML (rekord felvitel, törlés és módosítás), DCL és a \code{SELECT} utasítások áttekintése és használata. Beágyazott \code{SELECT} használata. A relációs algebra műveleteinek megvalósítása SQL-ben.
\\

\textbf{21.} &
Adatkezelés és adatbáziskezelés alapfogalmai, fileszervezési módszerek, B-fa index; adatbázis architektúra; A normalizálás szerepe a relációs adatbázis kezelésben. Redundancia oka és a megszüntetés lépései. Normálformák. Dekompozíció célja és szabályai, tételei.
\\

\hline

\textbf{22.} &
Számítógép-hálózatokhoz kötődő alapfogalmak és az ISO-OSI hivatkozási modell: topológia és méret szerinti osztályozásuk. Kapcsolástechnika szerinti osztályozásuk (vonalkapcsolás, üzenetkapcsolás, csomagkapcsolás és virtuális vonalkapcsolás). Az ISO-OSI hivatkozási modell szerkezete, rétegei és azok főbb funkciói.
\\

\textbf{23.} &
A számítógép-hálózatok ISO-OSI hivatkozási modell adatkapcsolati rétegének közeghozzáférési alrétege: A főbb csatorna-megosztási osztályok (statikus-dinamikus, dinamikus versengő-determinisztikus) bemutatása és azok összevetése. Az IEEE 802.3 és az Ethernet (CSMA/CD) keretformátum, MAC címek, támogatott közegek és sebességek bemutatása, működés duplex csatorna esetén. Az IEEE 802.11 WLAN általános felépítése, az Access Point (AP) főbb funkciói. Az alkalmazott közeghozzáférési módszer (CSMA/CA), Distributed Coordination Function (DCF) és Point Coordination Function (PCF) üzemmódok.
\\

\textbf{24.} &
A TCP/IP protokoll szövet és az Internet. Az Internet hivatkozási modell (DoD) és az ISO-OSI hivatkozási modell összevetése. A TCP/IP protokoll szövet főbb részei (ARP, RARP, IP, ICMP, TCP, UDP) és azok funkcióik. Az Internet címzés és címosztályok: IPv4 címosztályok, maszk, subnet, supernet, osztály nélküli címzés (CIDR -- Classless Inter-Domain Routing), és a változó alhálózat méretek (VLSM -- Variable Length Subnet Mask), címek kiosztása, lokális címek és a címfordítás (NAT). Az IPv6 címek, IPv6 cím típusok (unicast, multicast, anycast; link local, unique local, aggregatable global unicast address).
\\

\hline

\textbf{25.} &
Szoftvertechnológia és a szoftverfolyamat fogalma. A szoftverfolyamat fázisok részlet és bemutatása: szoftverspecifikáció, tervezés, implementáció, validáció és evolúció. Az evolúciós modell ismertetése.
\\

\textbf{26.} &
Szoftverkövetelmény fogalma és osztályozásai. Követelmények általános problémái. Követelmény feltárás és elemzés folyamata. Vezérlési stílusok.
\\

\textbf{27.} &
UML diagramok bemutatása: Use case és osztálydiagram. A szekvencia diagram. Komponens diagram.
\\

\hline

\end{tabular}

\end{document}
